% ---------------------------------------------------------------------------
% Pillar 2: Zero-Variance Fraction diagnostic --- Iter 42 early-trace ZVF
% identifiability of late failure mode (cross-pillar P1×P2 closure).
% Materialised from scripts/zvf_iter42.py outputs:
%   experiments/results/zvf_iter42_early_predicts.tsv
%   experiments/results/zvf_iter42_failure_audit.tsv
%   experiments/results/zvf_iter42_summary.tsv
% ---------------------------------------------------------------------------

\subsection{Early-Trace ZVF Identifiability of Late Failure Mode}
\label{sec:zvf-early-identify}
\S\ref{sec:zvf-iso-yield} closed the cross-library ZVF story by
asking what minimum group size $G$ a library needs to clear a target
contrastive yield $Y_{\mathrm{target}}$. This iter closes the
\emph{identifiability} loop by asking a different question: \emph{can
the failure mode of a training run be predicted from ZVF statistics
computed only over the first half of the trajectory?} The motivation
is operational: in production RL pipelines a researcher wants to
\emph{abort} a doomed run after a few hours, not after a few days.

The analysis pools the 60 (kind, method, seed) runs in
\texttt{experiments/results/zvf_dynamics.json}, where each run carries
per-step ZVF summaries (\texttt{first\_pass\_zvf05}, \texttt{auc\_above\_zvf05},
etc.) alongside the late-trace outcome (\texttt{last10\_avg},
\texttt{heldout\_acc}, \texttt{mean\_zvf}). We compute three
\emph{early-trace-only} features per run:
%
\begin{align*}
\mathrm{fp05\_frac} &:= \mathrm{first\_pass\_zvf05}/n_{\mathrm{steps}}
    \quad\text{(smaller $\Rightarrow$ ZVF climbs earlier)}\\
\mathrm{burden\_05} &:= \mathrm{auc\_above\_zvf05} \in [0,1]\\
\mathrm{burden\_07} &:= \mathrm{auc\_above\_zvf07} \in [0,1]
\end{align*}
%
We label each run with the deterministic failure taxonomy used by
iter~38: \texttt{collapse} if $\mathrm{last10\_avg}<0.05$;
\texttt{plateau} if $\mathrm{peak}<0.5$; \texttt{drift} if
$\mathrm{last10\_avg}<0.85\,\mathrm{peak}$; else \texttt{converged}.
On this pool the observed labels are exclusively \texttt{plateau}
($n{=}45$, all nine variance-mitigation libraries) and
\texttt{converged} ($n{=}15$, groupsize\_zvf\_sweep and
tinker\_prompt\_zvf); no run collapsed or drifted.

\paragraph{Leave-one-cluster-out kNN.}
A naive leave-one-row kNN over the 8 feature vector
(includes the full-trace \texttt{mean\_zvf}) trivially hits 60/60
accuracy because every seed in a given (kind, method) cluster shares
the same label. The harder question is whether the early features
\emph{generalise} across clusters. We therefore run a
\emph{leave-one-cluster-out} kNN ($k{=}3$, z-scored features) that
predicts every row's label using \emph{only the other 13 clusters}.
With only the 6 early-trace features (no \texttt{mean\_zvf} leakage)
the cluster-LOO accuracy is 58/60 ($0.967$), matching the full-feature
cluster-LOO exactly. Both of the two misclassified rows fall in the
\texttt{variance\_mitigation/grpo} cluster (early-feature cluster-LOO
accuracy $0.60$; all other clusters hit $1.000$).

\paragraph{Univariate AUC.}
The single most informative early feature is \texttt{burden\_05}: its
univariate AUC for the \texttt{converged} vs \texttt{plateau} binary
is $0.975$ (\texttt{burden\_07}: $0.978$). The earliest-crossing
features (\texttt{fp05\_frac}, \texttt{fp07\_frac}, \texttt{fp09\_frac})
by themselves have AUC $\approx 0$ on this pool --- the signal is in
the cumulative \emph{burden} above the threshold, not in the timing
of the first crossing. \texttt{zvf\_lag1} has AUC $0.105$ (the
autocorrelation structure does not separate the two classes). The
single-feature AUCs are summarised in
\texttt{experiments/results/zvf_iter42_summary.tsv}.

\paragraph{Cross-pillar P1$\times$P2 closure.}
Iter~41 (Pillar~1) showed that the saturation-model $R_{\max}$ is
identifiable from the first 60\% of the trace on 9/9 eligible anchors
(bootstrap 95\% CI on the prediction error contains 0 on 7/9). Iter~42
(Pillar~2) shows the dual statement for ZVF: on the
variance-mitigation + groupsize + tinker-prompt pool,
\emph{early-trace ZVF burden identifies the late failure mode without
any observation of the second half of the trace.} The two results
together license the operational claim: \emph{a researcher can
abort a doomed RL run within the first 30--50\% of the trace on the
basis of ZVF burden alone, without paying for the remaining compute.}

\paragraph{Why the burden, not the timing?}
The $\mathrm{AUC}_{\mathrm{fp05\_frac}} \approx 0$ vs
$\mathrm{AUC}_{\mathrm{burden\_05}} \approx 0.975$ asymmetry is
informative. The first-pass threshold captures \emph{when} ZVF first
spiked, but every cluster in this pool \emph{does} spike ZVF past
0.5 at some point (mean $\mathrm{fp05\_frac}$ is finite for every
cluster). What separates \texttt{converged} from \texttt{plateau} is
how long ZVF \emph{stays} above the threshold --- the cumulative
contrastive-loss burden --- not whether it ever crossed. This is
consistent with the iter-30 finding that the time-averaged
\texttt{trailing\_mean\_25} (a windowed burden statistic) is the
single most informative ZVF summary, and with iter-34's
phase-discriminator AUC above $0.97$.

\paragraph{Caveat.}
The 60-row pool carries only 2 of the 4 taxonomy classes. The
variance-mitigation libraries in this worktree do not exhibit
late-trace collapse (their \texttt{last10\_avg} hovers at $0.39$,
well above the $0.05$ threshold). The corresponding collapse runs
identified by the wider cross-experiment pool (Nemotron-120B GSM8K,
tool-use Qwen3-32B / Llama-8B-Instruct) carry $\mathrm{ZVF}{=}1.0$
throughout the entire trace and would be trivially separable by any
feature --- they would also be trivially separable by the i.i.d.\
baseline $\mathrm{ZVF}_{\mathrm{iid}}(p,G)$, which is what the
frontier synthesis flagged. The contribution of iter~42 is the
\emph{plateau-vs-converged} identifiability claim, not a re-derivation
of the trivial collapse case.

\paragraph{Reproduction.}
\texttt{scripts/zvf_iter42.py} (180 lines, stdlib only). Writes
\texttt{zvf_iter42_early_predicts.tsv} (60 rows $\times$ 21 cols),
\texttt{zvf_iter42_failure_audit.tsv} (14 clusters), and
\texttt{zvf_iter42_summary.tsv} (cluster-LOO, AUCs, ORs).
